\chapter{MÔ TẢ NGHIỆP VỤ}

\section{Quy trình phân tích -- thiết kế (tóm tắt)}

Báo cáo đi theo chuỗi chuẩn:

\begin{enumerate}[leftmargin=1.5cm]
    \item \textbf{Nghiệp vụ} (chương này) -- mô tả công việc doanh nghiệp.
    \item \textbf{Use case tổng quát} -- các nhóm chức năng hệ thống hỗ trợ.
    \item \textbf{Use case phân rã} -- mỗi hành động/mục tiêu của actor = 1 use case.
    \item \textbf{Đặc tả use case} -- actor, mục tiêu, luồng chính/thay thế.
    \item \textbf{Activity diagram} -- luồng hoạt động của từng use case (hoặc use case chính).
    \item \textbf{Sequence diagram} -- tương tác đối tượng theo thời gian.
    \item \textbf{Class diagram} -- lớp và quan hệ phục vụ các use case.
\end{enumerate}

\textit{Lưu ý}: Nghiệp vụ \textbf{không} đồng nhất với use case. Ví dụ nghiệp vụ ``Quản lý tài khoản'' gồm nhiều use case: thêm, sửa, xóa/khóa, tìm kiếm tài khoản.

\section{Tổng quan nghiệp vụ hệ thống}

Hệ thống quản lý kho phục vụ: xác thực người dùng; duy trì danh mục (tài khoản, NCC, khách hàng, hàng hóa, kho); nhập kho; xuất kho.

\begin{longtable}{|p{4.2cm}|p{8.8cm}|}
\hline
\textbf{Nghiệp vụ (chung)} & \textbf{Sẽ phân rã thành các use case (mục tiêu)} \\
\hline
\endfirsthead
\hline
\textbf{Nghiệp vụ (chung)} & \textbf{Sẽ phân rã thành các use case (mục tiêu)} \\
\hline
\endhead
Xác thực & Đăng nhập; Đăng xuất \\
\hline
Quản lý tài khoản & Thêm TK; Sửa TK; Xóa/khóa TK; Tìm TK \\
\hline
Quản lý nhà cung cấp & Thêm; Sửa; Xóa; Tìm NCC \\
\hline
Quản lý khách hàng & Thêm; Sửa; Xóa; Tìm KH \\
\hline
Quản lý hàng hóa & Thêm; Sửa; Ngừng; Tìm hàng \\
\hline
Quản lý kho & Thêm; Sửa; Tìm kho \\
\hline
Nhập kho & Lập phiếu nhập; Duyệt phiếu nhập \\
\hline
Xuất kho & Lập phiếu xuất; Duyệt phiếu xuất \\
\hline
\end{longtable}

\section{Nghiệp vụ 1: Xác thực người dùng}

Người dùng đăng nhập bằng tên đăng nhập và mật khẩu để sử dụng hệ thống; khi kết thúc thì đăng xuất (xóa phiên).

\textbf{Dữ liệu}: \texttt{NguoiDung}, \texttt{LichSuThaoTac}.

\section{Nghiệp vụ 2: Quản lý tài khoản}

Quản trị viên duy trì tài khoản và phân quyền (QTV / QL kho / NV kho). Gồm các hành động: thêm mới, sửa thông tin, khóa/xóa, xem/tìm danh sách.

\textbf{Dữ liệu}: \texttt{NguoiDung}.

\section{Nghiệp vụ 3: Quản lý nhà cung cấp}

Duy trì danh mục NCC phục vụ lập phiếu nhập: thêm, sửa, xóa, tìm kiếm.

\textbf{Dữ liệu}: \texttt{NhaCungCap}.

\section{Nghiệp vụ 4: Quản lý khách hàng}

Duy trì danh mục KH phục vụ lập phiếu xuất: thêm, sửa, xóa, tìm kiếm.

\textbf{Dữ liệu}: \texttt{KhachHang}.

\section{Nghiệp vụ 5: Quản lý hàng hóa}

Danh mục hàng (mã, tên, nhóm, ĐVT, giá): thêm, sửa, ngừng sử dụng, tìm kiếm.

\textbf{Dữ liệu}: \texttt{HangHoa}, \texttt{NhomHangHoa}, \texttt{DonViTinh}, \texttt{TonKho}.

\section{Nghiệp vụ 6: Quản lý kho}

Thông tin kho: thêm, sửa, tìm kiếm / xem tồn tại kho.

\textbf{Dữ liệu}: \texttt{Kho}, \texttt{TonKho}.

\section{Nghiệp vụ 7: Nhập kho}

Khi hàng về: nhân viên \textbf{lập phiếu nhập} (NCC, kho, dòng hàng) và gửi duyệt; quản lý \textbf{duyệt} thì hệ thống tăng tồn. Hai mục tiêu khác actor $\rightarrow$ hai use case (Chương 3).

\textbf{Dữ liệu}: \texttt{PhieuNhapKho}, \texttt{ChiTietPhieuNhap}, \texttt{TonKho}.

\section{Nghiệp vụ 8: Xuất kho}

Nhân viên \textbf{lập phiếu xuất} (KH, kho, dòng hàng, kiểm tra tồn), gửi duyệt; quản lý \textbf{duyệt} thì giảm tồn.

\textbf{Dữ liệu}: \texttt{PhieuXuatKho}, \texttt{ChiTietPhieuXuat}, \texttt{TonKho}.

\section{Kết luận chương}

Chương 2 mô tả \textbf{nghiệp vụ}. Bước tiếp theo (Chương 3): use case tổng quát $\rightarrow$ phân rã $\rightarrow$ đặc tả. Activity, sequence, class trình bày sau khi đã có use case.
